Esta entrega permitió implementar una metodología por etapas para planificar la operación diaria de buses eléctricos con expediciones y horarios establecidos, buscando minimizar el costo total. Se construyó un caso base para toda la red y se evaluó la propuesta completa en La Reina. El principal aprendizaje es que reducir las jornadas inicialmente asignadas no basta: la autonomía y la disponibilidad de puestos de carga condicionan la flota finalmente requerida, por lo que las decisiones de asignación y recarga deben coordinarse.

En toda la red, habilitar el \textit{interlining} con carga reactiva reduce el costo del caso base en un 5,2\%, aunque las restricciones energéticas absorben parte del beneficio inicial sobre la flota. En La Reina, la propuesta completa disminuye el costo diario de 594.602 a 424.300 USD y la flota de 2.045 a 1.383 buses, reducciones de 28,6\% y 32,4\%, respectivamente. La mejora proviene principalmente de reorganizar y reutilizar buses, y de programar las recargas para reducir reservas. El ahorro tarifario es secundario porque los puestos de carga ya se utilizan al máximo en los períodos más baratos. Esto sugiere que mejorar la coordinación operacional puede reducir la necesidad de flota adicional con la infraestructura existente, siempre que se corrijan las incompatibilidades de capacidad detectadas.

La principal fortaleza del enfoque es separar decisiones complejas en etapas evaluables, permitiendo identificar el origen de las mejoras. Respecto de su robustez, variar la discretización entre 5, 15 y 30 min modifica menos de un 0,5\% el costo de energía más reservas del modelo agregado. Sin embargo, esta estabilidad económica no demuestra factibilidad minuto a minuto. Asimismo, agregar holguras de recarga reduce reservas, pero aumenta la flota total y el costo: disponer de más tiempo no sustituye la coordinación con los puestos de carga. Queda pendiente evaluar la sensibilidad al consumo energético y a otros parámetros operacionales.

Los resultados siguen siendo preliminares. El programa detallado supera los 100 puestos disponibles durante 32 min, con un máximo de 115 buses, y persisten cadenas superiores a 24 h que afectan las reservas. Además, la propuesta completa solo se ha probado en un electroterminal. La brecha de optimalidad de 0,66\% corresponde al modelo detallado de carga y no garantiza la optimalidad global de la metodología. Su tiempo de resolución, cercano a una hora, también limita la extensión directa a instalaciones mayores.

Para la entrega final, las prioridades son corregir la capacidad simultánea y verificarla minuto a minuto, formular un ciclo diario estricto y retroalimentar la construcción de jornadas con la disponibilidad de puestos de carga. Luego, se debe extender la evaluación a toda la red y completar las sensibilidades pendientes. El modelo agregado ofrece una vía para reducir el tiempo de cómputo, pero requiere consolidar sus eventos de carga y validar el programa resultante. Estas mejoras permitirán determinar qué proporción del ahorro estimado se conserva en una operación factible, repetible y comparable con el caso base.
